% TOP_PROBLEM: {"id": 3100008, "title": "Does every thrackle have average degree at most 2", "category_id": 2, "category": {"id": 2, "name": "combinatorics", "display_name": "Combinatorics", "description": "Counting problems, graph theory, discrete structures.", "slug": "combinatorics", "order_index": 2, "created_at": "2026-07-31T15:26:25.671Z"}, "status": "partially_solved", "problem_number": "AMR-030-0008", "set_id": 13}
% FIRST_POSTED: 2026-10-02
% ENTRY_KIND: statement_only
% UnsolvedMath ID 3100008; source catalogue code AMR-030-0008
% !TeX program = lualatex
% Definition and sources only; this file is not an LLM solution attempt.
\documentclass[11pt,a4paper]{article}
\usepackage[margin=25mm]{geometry}
\usepackage{fontspec}
\setmainfont{lmroman10-regular.otf}[BoldFont=lmroman10-bold.otf,
 ItalicFont=lmroman10-italic.otf,BoldItalicFont=lmroman10-bolditalic.otf]
\usepackage{amsmath,amssymb,mathtools}
\usepackage{unicode-math}
\setmathfont{latinmodern-math.otf}
\AtBeginDocument{\let\setminus\smallsetminus}
\usepackage{xurl,hyperref}
\hypersetup{colorlinks=true,urlcolor=blue}
\setcounter{secnumdepth}{0}
\setlength{\parindent}{0pt}
\setlength{\parskip}{5pt}
\setlength{\emergencystretch}{3em}
\begin{document}
\section*{Does every thrackle have average degree at most 2}\label{3100008}
{\small UnsolvedMath ID 3100008 · AMR-030-0008 · Combinatorics · AMR Open Problem Lists}
\subsection{Definitions and mathematical statement}
A thrackle drawing of a finite simple graph $G=(V,E)$ in the plane
places distinct vertices at distinct points and represents every
edge by a simple Jordan arc joining its endpoints. An arc contains
no other vertex. Any two distinct edges meet exactly once: either
at their shared endpoint, when adjacent, or at one proper crossing
of their interiors, when nonadjacent. Tangencies and overlapping
subarcs are not counted as proper crossings and are excluded.
One may use general position, with no three edge interiors meeting
at a point.

Conway's thrackle conjecture is
\[
                           |E|\le |V|
\]
for every graph admitting such a drawing. The assertion concerns
the existence of a topological drawing, not whether a prescribed
placement of the vertices permits one.

Curved edges are essential to the scope of the question. Requiring
straight segments gives a more restrictive geometric problem and
does not resolve the conjecture for Jordan arcs. Adjacent edges
have already used their unique allowed meeting at their common
endpoint, so they may not cross again elsewhere. Isolated vertices
are allowed and contribute to $|V|$; deleting them only strengthens
a possible counterexample. A resolution must either establish the
inequality for every finite thrackle or exhibit a fully specified
eligible drawing with more edges than vertices.
\subsection{Short English statement}
Must every finite graph drawn as a planar thrackle have at most as many edges as vertices?
\subsection{Sources}
\begin{itemize}
\item[S1] ULAM.AI and the UnsolvedMath Contributors, supplied problems.json, record 3100008 (AMR-030-0008), AMR Open Problem Lists. Catalogue identity and source transcription; CC BY 4.0.
\url{https://huggingface.co/datasets/ulamai/UnsolvedMath}.
\item[S2] Cooper - Combinatorial Problems I Like (2020 snapshot), Geometry, source-order bullet 8. Original source indexed by AMR-030-0008.
\url{https://people.math.sc.edu/cooper/combprob.html}.
\end{itemize}
\end{document}
